%% ---------------------------------------------------------------------------------------
\section{Results: overfitting to 100\% training accuracy}
\label{sec:results}

Table~\ref{tab:main} is the main result. Every configuration reaches \textbf{100.0\% training
accuracy}; we verify this on every training sample after the tree is built, not from the training
loss. The columns ``root'' give the accuracy of the root network alone (what an ordinary small
network achieves), ``nodes (isol.)'' counts all networks in the tree, root included, and how many of
them are closed-form isolation nodes, and ``test fixed/broken'' attributes every test prediction that the
children changed: \emph{fixed} if the root was wrong and the tree is right, \emph{broken} if the root
was right and the tree is wrong.

\begin{table}[t]
\centering\small
\caption{Trees grown to 100\% training accuracy. Root = the root network alone. ``Test fixed/broken''
counts test images whose prediction the children changed from wrong to right / right to wrong.}
\label{tab:main}
\resizebox{\linewidth}{!}{\begin{tabular}{@{}lrrrrrrrr@{}}\toprule
 & \multicolumn{2}{c}{root acc. (\%)} & \multicolumn{2}{c}{tree acc. (\%)} & nodes & & & test fixed/\\
\cmidrule(lr){2-3}\cmidrule(lr){4-5}
configuration & train & test & train & test & (isol.) & params & depth & broken\\\midrule
\multicolumn{9}{@{}l}{\emph{CIFAR-10 -- corrector tree, root $128\to64\to10$, binary nodes $h=8$}}\\
pixels & 60.6 & 50.8 & 100.0 & \textbf{48.5} & 107 (6) & 1470.1k & 10 & 395/627\\
scratch CNN & 92.1 & 84.8 & 100.0 & \textbf{81.8} & 83 (2) & 144.0k & 6 & 109/406\\
ResNet-18 (ImageNet) & 84.2 & 74.5 & 100.0 & \textbf{69.8} & 67 (0) & 77.6k & 4 & 256/729\\
CNN + ResNet, per-node trunk selection & 92.1 & 84.8 & 100.0 & \textbf{82.2} & 35 (0) & 44.3k & 4 & 167/422\\
\midrule
\multicolumn{9}{@{}l}{\emph{Cats vs Dogs -- the note's binary tree, every node $128\to h\to1$}}\\
pixels & 71.5 & 63.2 & 100.0 & \textbf{59.8} & 38 (0) & 39.6k & 6 & 439/606\\
scratch CNN & 96.1 & 94.1 & 100.0 & \textbf{92.1} & 17 (1) & 17.4k & 4 & 54/155\\
ResNet-18 (ImageNet) & 97.9 & 93.4 & 100.0 & \textbf{93.2} & 5 (0) & 5.2k & 2 & 40/50\\
CNN + ResNet, per-node trunk selection & 97.9 & 93.4 & 100.0 & \textbf{93.3} & 5 (0) & 5.2k & 2 & 40/45\\
scratch CNN, single neurons ($h=0$, the note) & 95.8 & 93.6 & 100.0 & \textbf{93.6} & 3 (2) & 107.7k & 1 & 1/1\\
scratch CNN, $h=8$, LP (6)--(9) + $\ell_1$ & 96.2 & 94.0 & 100.0 & \textbf{92.1} & 17 (1) & 17.4k & 4 & 50/145\\
\bottomrule\end{tabular}
}
\end{table}

\begin{figure}[t]
\centering
\includegraphics[width=\linewidth]{tree_cd_cnn.pdf}
\caption{The Cats vs Dogs tree (CNN trunk, $h=8$), drawn like Fig.~4 of the note: each box is a
$128\!\to\!8\!\to\!1$ network trained on $n$ samples of which ``$+m$'' must make it fire; arrows carry
the closed-form connection weights \eqref{eq:closed}. A is always the ``missed'' child and B the
``false-alarm'' child. One child (red) could not make progress and became an isolation node with 6
caps. The tree classifies all 19\,989 training images correctly.}
\label{fig:treecd}
\end{figure}

\paragraph{Overfitting is always achieved, at very different costs.}
As with deep networks trained on random labels~\cite{zhang},
reaching 100\% is never the hard part. How many nodes interpolation needs is set by how far the root is from interpolating, i.e.\ by the
features. On Cats vs Dogs the ImageNet trunk leaves the root at 97.9\% training accuracy and the whole
tree has five nodes; with the scratch CNN it has 17, with pixels 38. On CIFAR-10 the root leaves 7.9k
(ResNet), 3.9k (CNN) or 19.7k (pixel) training errors and the corrector tree grows to
67, 83 and 107 nodes (the pixel tree reaches the depth limit and needs 1.47M parameters). In the paper's own setting -- every node a \emph{single} neuron
($h=0$) on CNN features -- the root neuron already reaches 95.8\%/93.6\%, but neither linear child
can solve its sub-problem: the balanced linear child $A$ catches 384 of the 404 missed dogs at the
price of 2\,076 false alarms. These ``inside-out'' sub-problems (positives surrounded by the
parent's correctly classified negatives) are exactly what made the MNIST tree of the first assignment
deep; here both children fall back to isolation nodes. With $h=8$ hidden units all but one child make
progress and the tree reaches depth 4 with a single isolation leaf.

\paragraph{The LP (6)--(9) and the redundant slacks.} Re-solving each node's output layer
together with $(w_A,w_B)$ as the note's LP (HiGHS, $\ell_1$ weight $10^{-3}$) gives \emph{exactly zero
total slack at every one of the 9 internal nodes} of the Cats vs Dogs tree -- the slacks are redundant,
as the note says -- and the same tree to within 0.06\% test accuracy (92.12\% vs 92.06\%). The LP
solutions are vertices with large connection weights ($w_A=232$, $w_B=-169$ at the root, versus
4.9 and $-3.8$ in closed form \eqref{eq:closed}), the badly scaled regime in which an LP solver can
report spurious non-zero slacks.

\paragraph{Interpolation lowers test accuracy, because the correctors generalise.}
Whenever the children are learned networks, the tree is \emph{worse} on test data than its own
root (by 0.1--4.7 points; Table~\ref{tab:main}), and the fixed/broken column shows why: each child is a small network trained on its
own few hundred or thousand positives, it learns a decision region rather than a list of points, and
on test images that region also covers images the root had right. Corrections break 1.3--4.4 test
predictions for every one they fix. This is the classical over-fitting picture, made visible node
by node.

\begin{table}[t]
\centering\small
\caption{Width of the binary nodes (CIFAR-10, ResNet trunk, corrector tree; root fixed at
$128\!\to\!64\!\to\!10$, 84.2\% train / 74.5\% test). ``Isolation only'' replaces every child by a
closed-form isolation node.}
\label{tab:width}
\resizebox{\linewidth}{!}{\begin{tabular}{@{}lrrrrrrrr@{}}\toprule
 & \multicolumn{2}{c}{root acc. (\%)} & \multicolumn{2}{c}{tree acc. (\%)} & nodes & & & test fixed/\\
\cmidrule(lr){2-3}\cmidrule(lr){4-5}
configuration & train & test & train & test & (isol.) & params & depth & broken\\\midrule
$h=0$ (linear $\to$ all isolation) & 84.2 & 74.5 & 100.0 & \textbf{74.6} & 11 (10) & 1027.5k & 1 & 10/0\\
$h=2$ & 84.2 & 74.5 & 100.0 & \textbf{73.3} & 93 (13) & 842.2k & 8 & 101/224\\
$h=8$ & 84.2 & 74.5 & 100.0 & \textbf{69.8} & 67 (0) & 77.6k & 4 & 256/729\\
$h=32$ & 84.2 & 74.5 & 100.0 & \textbf{70.7} & 18 (0) & 79.6k & 2 & 277/653\\
$h=128$ & 84.2 & 74.5 & 100.0 & \textbf{72.6} & 16 (0) & 258.5k & 2 & 284/475\\
isolation nodes only & 84.2 & 74.5 & 100.0 & \textbf{74.6} & 11 (10) & 1027.5k & 1 & 10/0\\
\bottomrule\end{tabular}
}
\end{table}

\paragraph{Benign versus harmful interpolation.} Table~\ref{tab:width} varies the width of the
correcting networks with the root held fixed. Linear children ($h=0$) never make progress, so the
tree consists of the root plus ten isolation nodes with 7\,902 caps -- one per training error -- and
its test accuracy is the root's: the caps changed only 10 of 10\,000 test predictions, all ten for
the better. The same holds when isolation is forced. Learned correctors are 13$\times$ smaller
(77.6k vs 1.03M parameters) but cost 1.2--4.7 points; the damage is largest for $h=8$--32 and
shrinks again for $h=128$, whose wider detectors fit their positives with tighter regions. This is
the sharpest form of the trade-off behind the whole assignment: \emph{memorising like a lookup table is
harmless but large; memorising with a small network is compact but generalises the noise.}

\begin{table}[t]
\centering\small
\caption{Multiclass extensions (CIFAR-10, ResNet trunk, $h=8$).}
\label{tab:archs}
\resizebox{\linewidth}{!}{\begin{tabular}{@{}lrrrrrrrr@{}}\toprule
 & \multicolumn{2}{c}{root acc. (\%)} & \multicolumn{2}{c}{tree acc. (\%)} & nodes & & & test fixed/\\
\cmidrule(lr){2-3}\cmidrule(lr){4-5}
configuration & train & test & train & test & (isol.) & params & depth & broken\\\midrule
one-vs-rest forest & 83.2 & 75.0 & 100.0 & \textbf{67.1} & 145 (0) & 150.9k & 5 & 467/1252\\
corrector tree & 84.2 & 74.5 & 100.0 & \textbf{69.8} & 67 (0) & 77.6k & 4 & 256/729\\
gated tree & 84.2 & 74.5 & 100.0 & \textbf{73.6} & 13 (1) & 918.3k & 5 & 27/118\\
\bottomrule\end{tabular}
}
\end{table}

\paragraph{Multiclass architectures (Table~\ref{tab:archs}).} The one-vs-rest forest is the most
expensive and the worst: ten independent roots make 10$\times$ more decisions, each binary root
errs on its own samples, and the forest needs 145 nodes and 151k parameters (67.1\% test). The
corrector tree shares one $K$-way root and grows detectors only for classes with errors (67 nodes,
77.6k). The gated tree is the most instructive: its re-classifier $R$ is easy (it relearns the 7\,902
errors of the root with only 31 mistakes), but its error detector $E$ must answer ``is the root wrong
here?'', which is as hard as the original problem -- the $h=8$ detector finds only 1\,060 of the 7\,902
errors and its missed-error child falls back to a 6\,839-cap isolation node. Detecting one's own errors
does not get easier by moving it into a separate node.

\begin{table}[t]
\centering\small
\caption{A trunk that has seen the data hides overfitting. CNN trunk trained on half $S_A$ of
CIFAR-10 (82.5\% test); trees grown on $S_A$ and on the unseen half $S_B$ (25k images each).}
\label{tab:half}
\resizebox{\linewidth}{!}{\begin{tabular}{@{}lrrrrrrrr@{}}\toprule
 & \multicolumn{2}{c}{root acc. (\%)} & \multicolumn{2}{c}{tree acc. (\%)} & nodes & & & test fixed/\\
\cmidrule(lr){2-3}\cmidrule(lr){4-5}
configuration & train & test & train & test & (isol.) & params & depth & broken\\\midrule
tree on $S_A$ (trunk's own training half) & 90.2 & 81.2 & 100.0 & \textbf{78.5} & 61 (0) & 71.4k & 6 & 149/427\\
tree on $S_B$ (half never seen by the trunk) & 85.0 & 81.7 & 100.0 & \textbf{76.3} & 95 (2) & 104.9k & 7 & 218/756\\
\bottomrule\end{tabular}
}
\end{table}

\paragraph{Where the memorisation lives (Table~\ref{tab:half}).} With features from a trunk that was
trained on $S_A$, the root fits $S_A$ to 90.2\% but the unseen $S_B$ only to 85.0\%, while both roots
test at $\approx$81.5\%. The 5-point gap is memorisation that already happened inside the trunk, and
the tree pays for it on $S_B$: 95 instead of 61 nodes and 47\% more parameters. When features come from
a trunk trained on the same images, a small tree is therefore \emph{not} evidence of an easy problem.
This is why the ImageNet trunk, which never saw the target images, is the cleanest place to study
interpolation in the tree itself, and why we use it for the double-descent experiments.

\paragraph{A set of trunks, chosen per node.} When every node may choose between the CNN and the
ResNet features (Table~\ref{tab:main}), the roots choose the trunk with the lower training error
(CNN on CIFAR-10, ResNet on Cats vs Dogs) but \emph{all ten} first-level detectors on CIFAR-10 choose the
ResNet features, and deeper nodes alternate (depth 2: 7 CNN / 7 ResNet; depth 3: 6/2). The root's
errors are exactly the images on which the CNN's features are ambiguous, so they are best separated in
a different feature space. The resulting tree is 2.4$\times$ smaller than the CNN-only tree (35 vs 83
nodes, 44k vs 144k parameters) and better on test data (82.2\% vs 81.8\%).
